% [原文 p.1]
\chapter{向量丛}
\section{基本定义}
\takeaway{本节界定向量空间族、向量丛、截面与同态，把“局部平凡”立为向量丛的唯一实质条件；此后的构造只在纤维上逐点进行。}

本讲义只展开复向量丛的理论。因此，除特别声明外，向量空间一律指复向量空间。不过，初等理论的大部分内容对实向量丛与辛向量丛同样成立。

设 $X$ 为拓扑空间。$X$ 上的向量空间族（family of vector spaces）由拓扑空间 $E$ 与以下两项组成：

\begin{enumerate}[label=(\roman*),leftmargin=2.4em,itemsep=2pt]
\item 连续映射 $p\colon E\to X$；
\item 对每个 $x\in X$，在
\[
E_x=p^{-1}(x)
\]
上给定的有限维向量空间结构，它与 $E$ 诱导到 $E_x$ 上的拓扑相容。
\end{enumerate}

映射 $p$ 称为投影，空间 $E$ 称为全空间，空间 $X$ 称为底空间；当 $x\in X$ 时，$E_x$ 称为 $x$ 上的纤维。

族 $p\colon E\to X$ 的截面（section）是指连续映射 $s\colon X\to E$，使对一切 $x\in X$ 有 $ps(x)=x$。

从族 $p\colon E\to X$ 到族 $q\colon F\to X$ 的同态映射（homomorphism）是指满足下述条件的连续映射 $\varphi\colon E\to F$：

\begin{enumerate}[label=(\roman*),leftmargin=2.4em,itemsep=2pt]
\item $q\varphi=p$；
\item 对每个 $x\in X$，$\varphi\colon E_x\to F_x$ 是向量空间之间的线性映射。
\end{enumerate}

% [原文 p.2]
当 $\varphi$ 为双射且 $\varphi^{-1}$ 连续时，称 $\varphi$ 为同构映射（isomorphism）。当 $E$ 与 $F$ 之间存在同构映射时，称 $E$ 与 $F$ 同构。

\begin{example*}{1}
  设 $V$ 为向量空间。令 $E=X\times V$，取 $p\colon E\to X$ 为向第一分量的投影。称 $E$ 为以 $V$ 为纤维的积族（product family）。\footnote{译注：积族 $X\times V$ 也记作 $\underline{V}$；特别地，当 $V=\C^n$ 时也记作 $\underline{n}$。（原日译本 p.2 译注）}若向量空间族 $F$ 同构于某个积族，则称 $F$ 为平凡（trivial）族。
\end{example*}

设 $Y$ 为 $X$ 的子空间。$E$ 为 $X$ 上的向量空间族，$p$ 为其投影，则 $p\colon p^{-1}(Y)\to Y$ 显然是 $Y$ 上的族。称此族为 $E$ 在 $Y$ 上的限制（restriction），记作 $E|_Y$。更一般地，取 $Y$ 为任意空间，$f\colon Y\to X$ 为连续映射，则诱导出的族 $f^*(p)\colon f^*(E)\to Y$ 定义如下：$f^*(E)$ 是 $Y\times E$ 的子空间，由满足 $f(y)=p(e)$ 的一切点 $(y,e)$ 组成，并配以显然的投影与向量空间结构。设 $g\colon Z\to Y$，则自然同构 $g^*f^*(E)\cong(fg)^*(E)$ 由把点 $(z,e)$（$z\in Z$，$e\in E$）逐个送到点 $(z,g(z),e)$ 给出。当 $f\colon Y\to X$ 为包含映射时，显然存在同构 $E|_Y\cong f^*(E)$，它把每个 $e\in E$ 送到对应的 $(p(e),e)$。

% [Tutor Note: 原书此处写“每个 $e\in E$”，按同构 $E|_Y\cong f^*(E)$ 的含义应读作 $e\in p^{-1}(Y)$；否则 $(p(e),e)$ 不在 $Y\times E$ 中。]

$X$ 上的向量空间族 $E$ 称为局部平凡（locally trivial），是指对每个 $x\in X$ 都存在邻域 $U$，使 $E|_U$ 平凡。局部平凡的族也称为向量丛（vector bundle）。平凡的族也称为平凡向量丛（平凡丛）。设 $f\colon Y\to X$，$E$ 是 $X$ 上的向量丛，则容易看出 $f^*(E)$ 是 $Y$ 上的向量丛。此时称 $f^*(E)$ 为诱导丛，即 $E$ 沿 $f$ 的拉回。

\begin{example*}{2}
  设 $V$ 为向量空间，$X$ 为与之相伴的射影空间。定义 $E\subset X\times V$ 为一切满足 $x\in X$、$v\in V$ 且 $v$ 位于 $x$ 所定的直线上的点 $(x,v)$ 所成的集合。$E$ 确为向量丛，其验证留给读者。
\end{example*}

% [原文 p.3]
当 $E$ 是 $X$ 上的向量丛时，注意 $\dim E_x$ 是 $X$ 上的局部常值函数，因而在 $X$ 的每个连通分支上是常值函数。当 $\dim E_x$ 在整个 $X$ 上为常数时，称 $E$ 具有维数，$E$ 的维数即各 $x$ 处公共的数 $\dim E_x$（注意：一般说来，如此定义的 $E$ 的维数，与 $E$ 作为拓扑空间的维数并不相同）。

向量丛局部平凡，因此向量丛的任一截面局部地由底空间上的向量值函数描述。当 $E$ 是向量丛时，把 $E$ 的一切截面组成的集合记作 $\Gamma(E)$。在任一空间上，取值于一个固定向量空间的函数全体自身构成向量空间，故可见 $\Gamma(E)$ 以自然的方式成为向量空间。

% [Tutor Note: 这里“$\Gamma(E)$ 是向量空间”只用到逐点线性结构，与截面空间上可赋的拓扑无关；原书此后也只用其线性结构，不涉及 $\Gamma(E)$ 的拓扑。]

设 $V, W$ 为向量空间，$E=X\times V$，$F=X\times W$ 为各自对应的积丛。此时任一同构映射 $\varphi\colon E\to F$ 由
\[
\varphi(x,v)=(x,\Phi(x)v)
\]
确定一个映射 $\Phi\colon X\to\Hom(V,W)$。再给 $\Hom(V,W)$ 以通常的拓扑，则 $\Phi$ 连续。反之，这样的连续映射 $\Phi\colon X\to\Hom(V,W)$ 确定同态映射 $\varphi\colon E\to F$（这一点在 $V, W$ 中取基 $\{e_i\}$、$\{f_i\}$ 时最容易看出。此时 $\Phi(x)$ 由满足下式的矩阵 $(\Phi(x)_{i,j})$ 表示：
\[
\Phi(x)e_i=\sum_j\Phi(x)_{i,j}f_j,
\]
$\varphi$ 与 $\Phi$ 的连续性和函数 $\Phi_{i,j}$ 的连续性等价）。

设 $\Iso(V,W)\subset\Hom(V,W)$ 为 $V$ 与 $W$ 之间一切同构映射组成的子空间。显然 $\Iso(V,W)$ 是 $\Hom(V,W)$ 的开集。此外，把 $T$ 映到逆映射 $T^{-1}$ 的映射给出连续映射 $\Iso(V,W)\to\Iso(W,V)$。设 $\varphi\colon E\to F$ 使每个 $x\in X$ 处的 $\varphi_x\colon E_x\to F_x$ 都是同构映射。这等价于 $\Phi(X)\subset\Iso(V,W)$。映射 $x\mapsto\Phi(x)^{-1}$ 定义 $\Psi\colon X\to\Iso(W,V)$，且连续。因此相应的映射 $\psi\colon F\to E$ 连续。于是，$\varphi\colon E\to F$ 为同构映射，当且仅当它是双射。再者，与此等价地，$\varphi$ 为同构映射，当且仅当各个 $\varphi_x$ 都是同构映射。此外，$\Iso(V,W)$ 是 $\Hom(V,W)$ 的开集，所以对任一同态映射 $\varphi$，使 $\varphi_x$ 为同构映射的点 $x\in X$ 的集合是 $X$ 的开子集。本质上，这些断言全是局部的，因此不仅对平凡族成立，对向量丛也成立。

% [Tutor Note: 此处“$\Iso(V,W)$ 显然为开集”用的是该集合由 $\det\neq0$ 刻画这一事实；$\dim V\neq\dim W$ 时 $\Iso(V,W)$ 为空集，论断仍成立。]

% [原文 p.4]
\begin{remark}{}
  $V$ 的有限维性在此前的论证中是本质的。若要考虑无限维向量丛，则必须区分 $\Hom(V,W)$ 上的各种算子拓扑来讨论。
\end{remark}

\section{向量丛上的运算}
\takeaway{只要函子 $T$ 在 $\Hom(V,W)$ 上连续，$T(E)$ 就获得自然拓扑；直和、张量积、$\Hom$、对偶丛与外幂都由此一次得到。}

直和、张量积等向量空间上的自然运算都可以扩张到向量丛上。唯一麻烦的问题是，如何给所得空间赋予拓扑。这里给出一条把运算从向量空间扩张到向量丛的一般途径，把一切运算的拓扑问题统一处理。

设 $T$ 为从有限维向量空间范畴到有限维向量空间范畴的函子。为简单起见，假定 $T$ 是单变元的共变函子。也就是说，对每个向量空间 $V$ 都有相应的向量空间 $T(V)$。若对一切 $V$ 与 $W$，映射 $T\colon\Hom(V,W)\to\Hom(T(V),T(W))$ 都连续，则称 $T$ 为连续函子（continuous functor）。

% [Tutor Note: 有限维向量空间之间的 $\Hom$ 上只有一种自然的向量空间拓扑，故这里说的连续性与拓扑的选取无关；到了无限维情形，$T$ 的连续性依赖所取的算子拓扑，这正是上一页所提醒的事。]

设 $E$ 为向量丛。定义集合 $T(E)$ 为和集
\[
\bigcup_{x\in X}T(E_x).
\]
于是对 $\varphi\colon E\to F$，通过映射 $T(\varphi_x)\colon T(E_x)\to T(F_x)$ 定义 $T(\varphi)\colon T(E)\to T(F)$。这里要证明的是，$T(E)$ 上可赋予自然的拓扑，且在此拓扑下 $T(\varphi)$ 连续。

先从 $E$ 为积丛的情形定义 $T(E)$。设 $E=X\times V$，把 $X\times T(V)$ 配以积拓扑，即作为 $T(E)$。设 $F=X\times W$，$\varphi\colon E\to F$ 为同态映射，$\Phi\colon X\to\Hom(V,W)$ 为与之对应的映射。由假设，$T\colon\Hom(V,W)\to\Hom(T(V),T(W))$ 连续，所以
% [原文 p.5]
$T\circ\Phi\colon X\to\Hom(T(V),T(W))$ 连续。因此 $T(\varphi)\colon X\times T(V)\to X\times T(W)$ 也连续。当 $\varphi$ 为同构映射时，$\varphi$ 连续且在每条纤维上是同构映射，故 $T(\varphi)$ 为同构映射。

其次，设 $E$ 平凡但未指定具体的积结构。选一同构映射 $\alpha\colon E\to X\times V$，赋予 $T(E)$ 拓扑，使 $T(\alpha)\colon T(E)\to X\times T(V)$ 为同胚。对另一任意的同构映射 $\beta\colon E\to X\times W$，令 $\varphi=\beta\circ\alpha^{-1}$，则 $T(\alpha)$ 与 $T(\beta)$ 在 $T(E)$ 上诱导相同的拓扑。这是因为 $T(\varphi)=T(\beta)\circ T(\alpha)^{-1}$ 是同胚。因此 $E$ 上的拓扑不依赖 $\alpha$ 的选取。此外，当 $Y\subset X$ 时，显然 $T(E)|_Y$ 上的拓扑与 $T(E|_Y)$ 上的拓扑相同。最后，当 $\varphi\colon E\to F$ 是平凡丛之间的同态映射时，$T(\varphi)\colon T(E)\to T(F)$ 连续，因而也是同态映射。

其次，设 $E$ 为向量丛。此时取 $U\subset X$ 使 $E|_U$ 平凡，则按上述方式 $T(E|_U)$ 上带有拓扑。$T(E)$ 的拓扑定义如下：子集 $V\subset T(E)$ 为开集，是指对每个使 $E|_U$ 平凡的 $U\subset X$，$V\cap(T(E)|_U)$ 都在 $T(E|_U)$ 中是开集。当 $Y\subset X$ 时，$T(E|_Y)$ 上的拓扑与 $T(E)|_Y$ 上的拓扑相同；又对任一同态映射 $\varphi\colon E\to F$，$T(\varphi)\colon T(E)\to T(F)$ 也是同态映射。这两件事都容易验证。

% [Tutor Note: 原书此处用 $V$ 表示 $T(E)$ 的任意子集，与上文的向量空间 $V$ 同名，只是局部的临时记法。]

设 $f\colon Y\to X$ 为连续映射，$E$ 为 $X$ 上的向量丛，则对任一连续函子 $T$，有自然同构
\[
(f^*\circ T)(E)\cong(T\circ f^*)(E).
\]
当 $T$ 是多变元函子、含共变成分与反变成分时，同样可以做。因此，对向量丛 $E, F$ 可以定义下列向量丛：

\begin{enumerate}[label=(\roman*),leftmargin=2.4em,itemsep=2pt]
\item $E\oplus F$（直和）；
\item $E\otimes F$（张量积）；
\item $\Hom(E,F)$；
\item $E^*$（$E$ 的对偶丛）；
\item $\lambda^i(E)$（$i$ 次外幂）。
\end{enumerate}

% [Tutor Note: 外幂的记号 $\lambda^i$ 在第 3 章将扩充为一整套 $\lambda$-运算 $\lambda_t(E)=\sum_i\lambda^i(E)t^i$，届时 $\lambda$ 同时充当“运算”与“外幂”两种读法，初读宜先把它固定为外幂。]

% [原文 p.6]
此外，还得到下列自然同构映射：

\begin{enumerate}[label=(\roman*),leftmargin=2.4em,itemsep=2pt]
\item $E\oplus F\cong F\oplus E$；
\item $E\otimes F\cong F\otimes E$；
\item $E\otimes(F'\oplus F'')\cong(E\otimes F')\oplus(E\otimes F'')$；
\item $\Hom(E,F)\cong E^*\otimes F$；
\item $\displaystyle\lambda^k(E\oplus F)\cong\bigoplus_{i+j=k}\bigl(\lambda^i(E)\otimes\lambda^j(F)\bigr)$。
\end{enumerate}

最后注意，$\Hom(E,F)$ 的截面与同态映射 $\varphi\colon E\to F$ 一一对应。因此，把 $\mathrm{HOM}(E,F)$ 定义为从 $E$ 到 $F$ 的一切同态映射组成的向量空间，并视 $\mathrm{HOM}(E,F)=\Gamma(\Hom(E,F))$。

\section{子丛与商丛}
\takeaway{单射、满射同态逐纤维判定；严格同态使核、像、余核仍是丛。子丛局部为直和项，商丛由此可定义；投影算子与 Hermite 度量是本节的技术收获。}

设 $E$ 为向量丛。$E$ 的子丛（sub-bundle）是指 $E$ 的子集，配以从 $E$ 限制而来的结构后本身成为向量丛。

当各个 $\varphi_x\colon F_x\to E_x$ 都是单射同态（或满射同态）时，称同态映射 $\varphi\colon F\to E$ 为单射同态（或满射同态）。注意，$\varphi\colon F\to E$ 为单射同态，当且仅当 $\varphi^*\colon E^*\to F^*$ 为满射同态。若 $F$ 是 $E$ 的子丛，$\varphi\colon F\to E$ 为包含映射，则 $\varphi$ 是单射同态。

\begin{lemma*}{1.3.1}
  若 $\varphi\colon F\to E$ 为单射同态，则 $\varphi(F)$ 是 $E$ 的子丛，且 $\varphi\colon F\to\varphi(F)$ 是同构映射。
\end{lemma*}

\begin{proof}
  $\varphi\colon F\to\varphi(F)$ 是双射，所以只要 $\varphi(F)$ 是子丛，$\varphi$ 就是同构映射。因此只需证明 $\varphi(F)$ 是子丛。

  问题局部，故考虑 $E$ 与 $F$ 为积丛的情形即足够。设 $E=X\times V$，取 $x\in X$。选 $W_x\subset V$ 为 $\varphi(F_x)$ 在 $V$ 中的补空间。$G=X\times W_x$ 是 $E$ 的子丛。

  用 $\theta(a\oplus b)=\varphi(a)+i(b)$ 定义 $\theta\colon F\oplus G\to E$，其中 $i\colon G\to E$ 为包含映射。按构造，$\theta_x$ 是同构映射。于是存在 $x$ 的开邻域 $U$，使 $\theta|_U$ 是同构映射。% [Tutor Note: 这里“存在开邻域 $U$ 使 $\theta|_U$ 为同构”用的是 p.4 末尾的结论：对平凡丛之间的同态，使某个纤维上为同构的点集是开集。原书只用“そうすると”一句带过。]
  由于 $F$ 是 $F\oplus G$ 的子丛，$\theta(F)=\varphi(F)$ 在 $U$ 上就是 $\theta(F\oplus G)=E$ 的子丛。
% [原文 p.7]
  于是，$\varphi(F)$ 是 $E$ 的子丛。
\end{proof}

注意，这个证明实际上给出了比引理更多的结论。也就是说，此论证对 $\varphi\colon F\to E$ 证明了：使 $\varphi_x$ 为单射同态的点集是开集。它还证明了局部地看，子丛是直和项。有了后一事实，就可以定义商丛。

\begin{definition}{1.3.1}
  设 $F$ 为 $E$ 的子丛。商丛 $E/F$ 定义为向量空间 $E_x/F_x$ 全体之和集配以商拓扑。
\end{definition}

% [Tutor Note: “配以商拓扑”这一句只交代了集合上的拓扑；$E/F$ 之所以是丛，靠的是上一段刚证明的“子丛局部是直和项”，由此每个 $E|_U/F|_U\cong U\times(V/W)$ 平凡。原书紧接着一句话补上了这一点。]

由于 $F$ 在 $E$ 中局部为直和项，$E/F$ 局部平凡，故它确实是向量丛。

设 $\varphi\colon F\to E$ 为任一同态映射，则函数 $\dim\ker\varphi_x$ 未必是常值函数，甚至未必是局部常值函数。

\begin{definition}{1.3.2}
  当 $\dim\ker\varphi_x$ 为局部常值函数时，称 $\varphi\colon F\to E$ 为严格同态映射（strict homomorphism）。
\end{definition}

\begin{proposition*}{1.3.2}
  若 $\varphi\colon F\to E$ 为严格同态映射，则下列 (i)–(iii) 成立。
  \begin{enumerate}[label=(\roman*),leftmargin=2.4em,itemsep=1pt]
  \item $\ker\varphi=\bigcup_x\ker\varphi_x$ 是 $F$ 的子丛；
  \item $\Image\varphi=\bigcup_x\Image\varphi_x$ 是 $E$ 的子丛；
  \item $\Coker\varphi=\bigcup_x\Coker\varphi_x$ 按商结构成为向量丛。
  \end{enumerate}
\end{proposition*}

\begin{proof}
  注意 (ii) 蕴含 (iii)。先证 (ii)。问题局部，故可设对某个 $V$ 有 $F=X\times V$。取 $x\in X$，选 $W_x\subset V$ 为 $\ker\varphi_x$ 在 $V$ 中的补空间。令 $G=X\times W_x$，则 $\varphi$ 与包含映射复合，诱导出同态映射 $\psi\colon G\to E$，使 $\psi_x$ 为单射。因此 $\psi$ 在 $x$ 的某个邻域 $U$ 上是单射。于是 $\psi(G)|_U$ 是 $E|_U$ 的子丛。另一方面 $\psi(G)\subset\varphi(F)$；又对一切 $y$ 有 $\dim\varphi(F_y)$ 为常数，而对一切 $y\in U$
\[
\dim\psi(G_y)=\dim\psi(G_x)=\dim\varphi(F_x)=\dim\varphi(F_y),
\]
% [原文 p.8]
故 $\psi(G)|_U=\varphi(F)|_U$ 成立。因此 $\varphi(F)$ 是 $E$ 的子丛。

最后必须证明 (i)。显然 $\varphi^*\colon E^*\to F^*$ 是严格同态映射。$F^*\to\Coker\varphi^*$ 是满射，故 $(\Coker\varphi^*)^*\to F^{**}$ 是单射。另一方面，对每个 $x$ 有自然的交换图
\[
\begin{array}{ccc}
\ker\varphi_x & \longrightarrow & F_x\\
\downarrow & & \downarrow\\
(\Coker\varphi_x^*)^* & \longrightarrow & F_x^{**}
\end{array}
\]
其中竖直方向的射都是同构映射。因此 $\ker\varphi\cong(\Coker\varphi^*)^*$，由引理 1.3.1，它是 $F$ 的子丛。
\end{proof}

% [Tutor Note: 原书这里的交换图只画了方块，未标出两条水平射各自是什么；上水平射是 $\ker\varphi_x\subset F_x$ 的包含，下水平射是把 $f\in F_x^{**}$ 限制到 $\Coker\varphi_x^*$ 上所得的线性泛函。]

这个证明同样给出了命题所述之外的结论。此论证证明了：对任意 $x\in X$，都存在 $x$ 的邻域 $U$，使一切 $y\in U$ 满足 $\dim\varphi_x(F_x)\le\dim\varphi_y(F_y)$。因此 $\operatorname{rank}\varphi_x$ 是 $x$ 的上半连续函数。

\begin{definition}{1.3.3}
  投影算子 $P\colon E\to E$ 是指满足 $P^2=P$ 的同态映射。
\end{definition}

% [Tutor Note: 这里把幂等同态称作“投影算子”，是纯代数的叫法；它未必是到子丛的正交投影，因为此时还没有度量。]

由于 $\operatorname{rank}P_x+\operatorname{rank}(1-P_x)=\dim E_x$，而 $\operatorname{rank}P_x$ 与 $\operatorname{rank}(1-P_x)$ 都是 $x$ 的上半连续函数，故它们都是局部常值函数。因此 $P$ 与 $1-P$ 是严格同态映射。由于 $\ker P=(1-P)E$，$E$ 是向量丛 $PE$ 与 $(1-P)E$ 的直和。因此，任一投影算子 $P\colon E\to E$ 都确定直和分解 $E=(PE)\oplus((1-P)E)$。

% [Tutor Note: “两个上半连续函数之和为常值函数，故两者都局部常值”这一步是原书一带而过的关键：上半连续加下半连续才得连续，此处靠的是 $\operatorname{rank}P_x+\operatorname{rank}(1-P_x)\equiv\dim E_x$ 常数这一恒等式。]

下面考虑向量丛上的度量。函子 $\mathrm{Herm}$ 把向量空间 $V$ 对应到由 $V$ 上一切 Hermite 形式组成的向量空间 $\mathrm{Herm}(V)$。用 1.2 节的方法，可以对向量丛 $E$ 定义向量丛 $\mathrm{Herm}(E)$。

\begin{definition}{1.3.4}
  向量丛 $E$ 上的度量（metric）是指截面 $h\colon X\to\mathrm{Herm}(E)$，使对每个 $x\in X$，$h(x)$ 都是正定的。选定了度量的向量丛称为 Hermite 丛。
\end{definition}

% CHUNK ch01a DONE
